\documentclass[12pt,a4paper]{article}

% Core Packages
\usepackage[utf8]{inputenc}
\usepackage[T1]{fontenc}
\usepackage{geometry}
\usepackage{graphicx}
\usepackage{hyperref}
\usepackage{titlesec}
\usepackage{enumitem}
\usepackage{booktabs}
\usepackage{array}
\usepackage{xcolor}
\usepackage{fancyhdr}
\usepackage{setspace}
\usepackage{parskip}
\usepackage{caption}
\usepackage{float}
\usepackage{amsmath}
\usepackage{amssymb}
\usepackage{algorithm}
\usepackage{algpseudocode}

% Page geometry
\geometry{margin=1in}

% Colors
\definecolor{indigoblue}{RGB}{0, 51, 102}
\definecolor{sectioncolor}{RGB}{0, 76, 153}

% Hyperref setup
\hypersetup{
    colorlinks=true,
    linkcolor=indigoblue,
    urlcolor=indigoblue,
    citecolor=indigoblue
}

% Section formatting
\titleformat{\section}{\Large\bfseries\color{sectioncolor}}{\thesection}{1em}{}
\titleformat{\subsection}{\large\bfseries\color{indigoblue}}{\thesubsection}{1em}{}

% Header/Footer
\pagestyle{fancy}
\fancyhf{}
\fancyhead[L]{\small COL761: Data Mining}
\fancyhead[R]{\small Assignment 1 - Part 3}
\fancyfoot[C]{\thepage}
\renewcommand{\headrulewidth}{0.4pt}

\onehalfspacing

\begin{document}

% ==============================================================================
% COVER PAGE
% ==============================================================================
\begin{titlepage}
    \centering
    \vspace*{1cm}
    {\Huge\bfseries\color{indigoblue} COL761 Data Mining}\\[0.5cm]
    {\LARGE Assignment 1}\\[0.3cm]
    {\Large\textbf{Part 3: Graph Indexing}}\\[2cm]
    
    {\LARGE\bfseries Team VindyCS5}\\[0.5cm]
    
    \begin{tabular}{ll}
        \toprule
        \textbf{Name} & \textbf{Entry Number} \\
        \midrule
        Abhiram Dharme & 2022CS51139 \\
        Raj Chandak & 2022CS51643 \\
        Purushottam Sharma & 2022CS51635 \\
        \bottomrule
    \end{tabular}\\[1cm]
    
    \rule{\textwidth}{0.5pt}\\[0.3cm]
    {\large Submitted: February 2, 2026}
    \vfill
\end{titlepage}

\tableofcontents
\newpage

%==============================================================================
% PROBLEM STATEMENT
%==============================================================================
\section{Problem Statement}
\textit{[10 (report) + 35 (competitive scoring) marks]}

Given a database of graphs $\mathcal{D} = \{g_1, g_2, \ldots, g_n\}$ and a query graph $q$, the goal is to efficiently identify a \textbf{candidate set} $C_q \subseteq \mathcal{D}$ containing all graphs that potentially contain $q$ as a subgraph.

The performance metric is:
\[
s_q = \frac{|R_q|}{|C_q|}
\]
where $|R_q|$ is the number of true positives (graphs actually containing $q$) and $|C_q|$ is the candidate set size. \textbf{Larger scores indicate better pruning.}

\subsection{Key Insight}

The necessary condition for subgraph isomorphism states: if $q \subseteq g_i$, then every substructure (fragment) present in $q$ must also be present in $g_i$. By identifying \textbf{discriminative subgraph fragments}, we can efficiently prune the candidate set.

%==============================================================================
% APPROACH
%==============================================================================
\section{Approach: gIndex-Inspired Feature Mining}

Our approach is inspired by the \textbf{gIndex} framework \cite{yan2004gindex}, which identifies discriminative subgraph fragments to construct a feature-based index. We mine multiple types of substructures and select the most discriminative ones based on variance scoring.

\subsection{Feature Types}

We mine six categories of subgraph features:

\begin{table}[H]
\centering
\begin{tabular}{llcc}
\toprule
\textbf{Feature Type} & \textbf{Description} & \textbf{Max Selected} & \textbf{Selection Method} \\
\midrule
Node Labels & Individual vertex labels & 20 & Frequency (top-k) \\
Edge Labels & Canonical edge triples $(l_u, e, l_v)$ & 30 & Frequency (top-k) \\
Path-3 & 3-node paths (2 edges) & 90 & Variance \\
Path-4 & 4-node paths (3 edges) & 160 & Variance \\
Cycle-5 & 5-node cycles & 40 & Variance \\
Cycle-6 & 6-node cycles & 50 & Variance \\
\bottomrule
\end{tabular}
\caption{Feature Types and Selection Counts}
\end{table}

\subsection{Canonical Representation}

Each feature type uses a \textbf{canonical form} to avoid redundant features:

\begin{itemize}
    \item \textbf{Edge}: Ordered as $(l_u, e, l_v)$ where $l_u \leq l_v$
    \item \textbf{Path-3}: Represented as $(l_a, e_1, l_b, e_2, l_c)$ with lexicographic ordering
    \item \textbf{Path-4}: Forward/backward ordering, choosing the lexicographically smaller
    \item \textbf{Cycles}: Sorted node labels + sorted edge labels (order-independent)
\end{itemize}

%==============================================================================
% FEATURE SELECTION
%==============================================================================
\section{Feature Selection: Variance-Based Scoring}

\subsection{Motivation}

Not all mined substructures are equally useful for pruning. Features that appear in \textbf{too many} or \textbf{too few} graphs provide little discriminative power:

\begin{itemize}
    \item Features in 99\% of graphs: Almost never prune anything
    \item Features in 1\% of graphs: Rarely present in queries
\end{itemize}

\subsection{Variance Score}

We use the \textbf{Bernoulli variance} as the discriminative score:

\[
\text{score}(f) = p(1-p) \quad \text{where} \quad p = \frac{\text{count}(f)}{N}
\]

This score is maximized when $p = 0.5$ (feature present in exactly half the graphs), providing maximum discriminative power.

\subsection{Selection Algorithm}

\begin{algorithm}[H]
\caption{Variance-Based Feature Selection}
\begin{algorithmic}[1]
\Require Frequency map $\texttt{freq}$, database size $N$, budget $k$
\Ensure Selected features $S$
\State $\texttt{scored} \gets []$
\For{each feature $f$ with count $c$ in $\texttt{freq}$}
    \State $p \gets c / N$
    \State $\text{score} \gets p \times (1 - p)$
    \If{$\text{score} > 0$}
        \State Append $(\text{score}, f)$ to $\texttt{scored}$
    \EndIf
\EndFor
\State Sort $\texttt{scored}$ in descending order
\State \Return top $k$ features from $\texttt{scored}$
\end{algorithmic}
\end{algorithm}

%==============================================================================
% FEATURE EXTRACTION
%==============================================================================
\section{Feature Extraction via Subgraph Isomorphism}

\subsection{Convert Phase}

Once discriminative features are identified, we convert each graph (database or query) into a \textbf{binary feature vector}:

\[
v_g = (b_1, b_2, \ldots, b_k) \in \{0,1\}^k
\]

where $b_j = 1$ if and only if feature graph $f_j$ is subgraph-isomorphic to $g$.

\subsection{Implementation}

We use \textbf{rustworkx} (Rust-based graph library) for efficient subgraph isomorphism testing:

\begin{verbatim}
is_present = rx.is_subgraph_isomorphic(
    graph, feature,
    node_matcher=lambda x, y: x['label'] == y['label'],
    edge_matcher=lambda x, y: x['label'] == y['label'],
    induced=False
)
\end{verbatim}

This provides significant speedup over pure Python implementations.

%==============================================================================
% CANDIDATE GENERATION
%==============================================================================
\section{Candidate Set Generation}

\subsection{Necessary Condition}

For a database graph $g_i$ to contain query $q$ as a subgraph, every feature present in $q$ must also be present in $g_i$:

\[
g_i \in C_q \iff v_q \leq v_{g_i} \quad \text{(component-wise)}
\]

\subsection{Algorithm}

\begin{algorithm}[H]
\caption{Candidate Set Generation}
\begin{algorithmic}[1]
\Require Database features $V_D$, query features $v_q$
\Ensure Candidate set $C_q$
\State $C_q \gets \emptyset$
\For{each database graph $g_i$ with features $v_{g_i}$}
    \If{$v_q \leq v_{g_i}$ (all query features present)}
        \State Add $g_i$ to $C_q$
    \EndIf
\EndFor
\State \Return $C_q$
\end{algorithmic}
\end{algorithm}

\subsection{Optimization}

The component-wise comparison $v_q \leq v_{g_i}$ is implemented using NumPy's vectorized operations:

\begin{verbatim}
if np.all(query_vector <= db_vector):
    candidates.append(graph_id)
\end{verbatim}

This allows batch processing of all database graphs efficiently.

%==============================================================================
% IMPLEMENTATION DETAILS
%==============================================================================
\section{Implementation Details}

\subsection{Pipeline Overview}

\begin{enumerate}
    \item \textbf{identify.sh}: Mine discriminative subgraphs from database
    \begin{itemize}
        \item Parse database graphs
        \item Mine node, edge, path, and cycle features
        \item Select top-k by frequency (nodes/edges) or variance (paths/cycles)
        \item Save as pickled rustworkx graphs
    \end{itemize}
    
    \item \textbf{convert.sh}: Convert graphs to feature vectors
    \begin{itemize}
        \item Load discriminative subgraphs
        \item For each graph, test subgraph isomorphism against each feature
        \item Output binary NumPy matrix
    \end{itemize}
    
    \item \textbf{generate\_candidates.sh}: Generate candidate sets
    \begin{itemize}
        \item Load database and query feature matrices
        \item For each query, find databases where $v_q \leq v_{g_i}$
        \item Output in required format
    \end{itemize}
\end{enumerate}

\subsection{Libraries Used}

\begin{table}[H]
\centering
\begin{tabular}{ll}
\toprule
\textbf{Library} & \textbf{Purpose} \\
\midrule
rustworkx & Graph representation and subgraph isomorphism \\
NumPy & Feature matrix operations \\
pickle & Serialization of feature graphs \\
\bottomrule
\end{tabular}
\caption{External Libraries}
\end{table}

%==============================================================================
% DESIGN DECISIONS
%==============================================================================
\section{Design Decisions and Rationale}

\subsection{Why Multiple Feature Types?}

Different query graphs have different structural characteristics. By including multiple feature types:

\begin{itemize}
    \item \textbf{Node/Edge labels}: Prune based on basic structural requirements
    \item \textbf{Path-3/Path-4}: Capture local connectivity patterns
    \item \textbf{Cycle-5/Cycle-6}: Capture cyclic structures common in molecular graphs
\end{itemize}

\subsection{Why Variance Scoring?}

Variance scoring ($p(1-p)$) naturally balances:
\begin{itemize}
    \item Features that are too common (low pruning power)
    \item Features that are too rare (unlikely to appear in queries)
\end{itemize}

Features with $p \approx 0.5$ provide maximum information for distinguishing graphs.

\subsection{Why rustworkx?}

rustworkx provides:
\begin{itemize}
    \item Rust-level performance for subgraph isomorphism
    \item Python API for ease of use
    \item Significant speedup over NetworkX (10-100$\times$)
\end{itemize}

%==============================================================================
% CONCLUSION
%==============================================================================
\section{Conclusion}

Our gIndex-inspired approach combines:
\begin{enumerate}
    \item \textbf{Diverse feature mining}: Node labels, edges, paths, and cycles
    \item \textbf{Variance-based selection}: Maximizing discriminative power
    \item \textbf{Efficient implementation}: rustworkx for fast isomorphism testing
\end{enumerate}

This approach generates compact candidate sets by exploiting the necessary condition for subgraph isomorphism across multiple structural feature types.

%==============================================================================
% REFERENCES
%==============================================================================
\newpage
\begin{thebibliography}{9}

\bibitem{yan2004gindex}
Yan, X., Yu, P. S., \& Han, J. (2004).
\textit{Graph indexing: A frequent structure-based approach}.
ACM SIGMOD, 335--346.

\bibitem{cheng2007fg}
Cheng, J., Ke, Y., Ng, W., \& Lu, A. (2007).
\textit{Fg-index: towards verification-free query processing on graph databases}.
ACM SIGMOD, 857--872.

\bibitem{shang2008gsting}
Shang, H., Zhang, Y., Lin, X., \& Yu, J. X. (2008).
\textit{Taming verification hardness: An efficient algorithm for testing subgraph isomorphism}.
VLDB, 1--12.

\end{thebibliography}

\end{document}
